\icode{N15J-MP-MU-L2-261004-005}
\nceeID[6]{2014天津}{8}
（多选）
一束由两种频率不同的单色光组成的复色光从空气射入玻璃三棱镜后，出射光分成\(a\)、\(b\)两束，如图所示，则\(a\)、\(b\)两束光\paren[]

\begin{tikzpicture}[scale=.8,decoration={markings,mark=at position 0.6 with {\arrow{>}}}]
	\draw (0,0)--(70:3)--(110:3)--(0,0);
	\draw[postaction={decorate},purple] (-2,2)--(110:1.5);
	\draw[postaction={decorate},purple] (110:1.5)--(70:1.3);
	\draw[postaction={decorate},purple] (110:1.5)--(70:1.45);
	\draw[postaction={decorate},purple] (70:1.45)--(1.7,2)node[right]{\(a\)};
	\draw[postaction={decorate},purple] (70:1.3)--(1.7,1.6)node[right]{\(b\)};
\end{tikzpicture}
\begin{choices}
	\item 垂直穿过同一块平板玻璃，\(a\)光所用的时间比\(b\)光长
	\item 从同种介质射入真空发生全反射时，\(a\)光临界角比\(b\)光的小
	\item 分别通过同一双缝干涉装置，\(b\)光形成的相邻条纹间距小
	\item 若照射同一金属都能发生光电效应，\(b\)光照射时逸出的光电子最大初动能大
\end{choices}
\begin{answer}
	AB\qquad 提示：\(n = \frac{c}{v}\)折射率越大速度越小；\(\sin C = \frac{1}{n}\)；\(\Delta x = \frac{L}{d}\lambda\)
\end{answer}